%%%%%%%%%%  scan idx 227  |  folio 214  %%%%%%%%%%
\SGAPrintPage{214}%
Identifying $u$, by (3.1.4), with a continuous map from $S$ to $\Diff^n(E,F)$, we see that, for $U$ open in $S$, the morphism
\[
  \pr_{1*}(u)(U):\Gamma(X|U,E)\longrightarrow \Gamma(X|U,F)
\]
is canonically identified with the linear map
\[
  C(U,C^\infty(V,E))\longrightarrow C(U,C^\infty(V,F))
\]
deduced from
\begin{equation}\tag{3.1.3}
  \phi u:U\longrightarrow \Lop(C^\infty(V,E),C^\infty(V,F)).
\end{equation}
Suppose $S$ is locally compact. Then (3.1.1) the morphism
\[
  \pr_{1*}(u):\pr_{1*}\pr_2^*(E)
  \longrightarrow \pr_{1*}\pr_2^*(F)
\]
is the morphism of sheaves associated to the morphism of presheaves
\[
  u_\infty:\underline C'(S,C^\infty(V,E))
  \longrightarrow \underline C'(S,C^\infty(V,F))
\]
defined by $\phi u$. On the other hand, for $r\in\underline{\underline{Z}}$,
\[
  \phi_r u:S\longrightarrow \Lop(C^r_2(V,E),C^{r-n}_2(V,F))
\]
defines a morphism of presheaves
\[
  u_r:\underline C'(S,C^r_2(V,E))
  \longrightarrow \underline C'(S,C^{r-n}_2(V,F)).
\]
By (3.1.3 (ii)), the $u_r$ form a projective system whose limit is $u_\infty$. It will be convenient to interpret $\phi_r u$ as a morphism of trivial Hilbert bundles (3.0.5)
\[
  S\times C^r_2(V,E)\longrightarrow S\times C^{r-n}_2(V,F),
\]
%%%%%%%%%%  scan idx 228  |  folio 215  %%%%%%%%%%
\SGAPrintPage{215}%
inducing on the sheaves of sections the morphism of sheaves associated to $u_r$.

\sgaStableTarget{sga6.e02.sec.3.2}\subsubsection*{3.2. Definition of the index}
\sgaStableTarget{sga6.e02.thm.3-2-1}\textbf{Theorem 3.2.1.} \uline{Let} $S$ \uline{be a locally compact space,} $p:X\to S$ \uline{a proper mixed variety,} $E^\bullet$ \uline{an elliptic complex of order} $\underline{n}$ \uline{of complex vector bundles on} $X$ (1.2.5)\uline{. Then the complex} $\underline{\underline{R}}p_*(E^\bullet)$ \uline{(which, by} (2.5)\uline{, is isomorphic, to} $p_*(E^\bullet)$ \uline{in the derived category of that of} $\underline{O}_S\otimes\underline{\underline{C}}$\uline{-Modules) is perfect.}

\uline{Proof.} The question being local on $S$, we may assume (2.9) that $X=S\times V$, where $V$ is a compact variety, and $p=\pr_1$. Moreover, we may assume (2.10) that, for every $i$, $E^i=\pr_2^*(F^i)$, where $F^i$ is a complex vector bundle of finite rank on $V$. By (3.1.5), the complex $p_*(E)$ is associated to the complex of presheaves $\bigl(\underline C'(S,C^\infty(V,F^i)),d_\infty^i\bigr)$ (where the $d^i:E^i\longrightarrow E^{i+1}$ are the given differential operators), which we shall denote simply by $\underline C'(S,C^\infty(V,F^\bullet))$. Put, for each $i\in\underline{\underline{Z}}$,
\[
  m_i=\sum_{j<i} n_j
  \qquad (\text{we agree that } n_i=0 \text{ when } E^i=0).
\]
For each $r\in\underline{\underline{Z}}$, we have (3.1.5) a complex of presheaves $\bigl(\underline C'(S,C_2^{r-m_i}(V,F^i)),d_{r-m_i}^i\bigr)$, which we shall denote by $\underline C'(S,C_2^r(V,F^\bullet))$.
The $\underline C'(S,C_2^r(V,F^\bullet))$ form a projective system, and we have (3.1.5):
\sgaStableTarget{sga6.e02.eq.3-2-1-1}\begin{equation}\tag{3.2.1.1}
  \underline C'(S,C^\infty(V,F^\bullet))=\varprojlim_r \underline C'(S,C_2^r(V,F^\bullet)).
\end{equation}
Finally, we shall denote by $F_r^\bullet$ the complex of trivial Hilbert bundles over $S$ defined by $\phi_{r-m_i}d^i$, and by $\underline F_r^\bullet$ the corresponding complex of sheaves of sections, which is precisely (3.1.5) the complex of sheaves associated to $\underline C'(S,C_2^r(V,F^\bullet))$. With this notation, to prove the theorem it suffices to prove the following assertions:

(a) \uline{For} $r\in\underline{\underline{Z}}$\uline{, the canonical morphism}
%%%%%%%%%%  scan idx 229  |  folio 216  %%%%%%%%%%
\SGAPrintPage{216}%
\[
  \underline C'(S,C^\infty(V,F^\bullet))
  \longrightarrow \underline C'(S,C_2^r(V,F^\bullet))
\]
\uline{is a
quasi-isomorphism of complexes of presheaves, i.e.\ induces isomorphisms on the cohomology presheaves.}

(b) \uline{For} $r\in\underline{\underline{Z}}$\uline{, the complex} $\underline F_r^\bullet$ \uline{is perfect.}

\uline{Proof of} (b). The ellipticity of $E^\bullet$ implies ([16]~XI Th.6\SGAAuthorFootnote{1}{which readily generalizes to elliptic complexes.}) that, for every $s\in S$, the differential of $(F_r^\bullet)_s$ has closed image and the $H^i((F_r^\bullet)_s)$ are finite-dimensional vector spaces. In other words, to use the terminology of [11], $F_r^\bullet$ is a direct quasi-acyclic complex of Hilbert bundles. Consequently, by ([12]~2.3), $F_r^\bullet$ is locally homotopically equivalent to a bounded complex of vector bundles of finite rank. Since the \textquotedblleft sheaf of sections\textquotedblright{} functor is additive (3.0.5), $\underline F_r^\bullet$ is locally homotopically equivalent to a bounded complex of locally free sheaves of finite type, and hence is perfect.

For the proof of (a), we shall need the following intermediate assertion:

(a') \uline{For} $p\geq q$\uline{, the canonical morphism}
\[
  \underline F_p^\bullet\longrightarrow \underline F_q^\bullet
\]
\uline{is a quasi-isomorphism.}

\uline{Proof of} (a'). Denote by $M_{pq}^\bullet$ the complex of Hilbert bundles that is the cone of the morphism $F_p^\bullet\longrightarrow F_q^\bullet$. The ellipticity of $E$ implies that, for every $s\in S$, the morphism $(F_p^\bullet)_s\longrightarrow (F_q^\bullet)_s$ is a quasi-isomorphism ([16]~XI Cor.~2 of Th.~6). In other words, $M_{pq}^\bullet$ is a direct acyclic complex. It follows ([11]~2.3) that $M_{pq}^\bullet$ is locally homotopically trivial, and hence that $\underline M_{pq}^\bullet$, which is precisely the cone of $\underline F_p^\bullet\longrightarrow\underline F_q^\bullet$, is acyclic, q.e.d.

\uline{Proof of} (a). We must show that, for every relatively compact open subset $U$ of $S$, the morphism
%%%%%%%%%%  scan idx 230  |  folio 217  %%%%%%%%%%
\SGAPrintPage{217}%
\[
  C(\overline U,C^\infty(V,F^\bullet))
  \longrightarrow C(\overline U,C_2^r(V,F^\bullet))
\]
is a
quasi-isomorphism. Replacing $S$ by $\overline U$ (and $E^\bullet$ by the complex induced over $\overline U$), we are reduced to showing that
\[
  C(S,C^\infty(V,F^\bullet))
  \longrightarrow C(S,C_2^r(V,F^\bullet))
\]
is a quasi-isomorphism, with $S$ assumed \uline{compact}. By (3.2.1.1), we have
\[
  C(S,C^\infty(V,F^\bullet))=
  \varprojlim_r C(S,C_2^r(V,F^\bullet))
\]
For each $i$, $C(S,C_2^r(V,F^i))$ is a Banach space (for the topology of uniform convergence), and the transition morphisms
\[
  C(S,C_2^p(V,F^i))\longrightarrow C(S,C_2^q(V,F^i))
\]
are, by (3.0.4) and (3.1.2), continuous linear maps with dense image. On the other hand, applying to the quasi-isomorphisms (a') $\underline F_p^\bullet\longrightarrow\underline F_q^\bullet$ the functor $\Gamma_S$, which is exact since $\underline{O}_S$ is soft, we see that the transition morphisms
\[
  C(S,C_2^p(V,F^\bullet))\longrightarrow C(S,C_2^q(V,F^\bullet))
\]
are quasi-isomorphisms. Assertion (a) then follows from (EGA $0_{\mathrm{III}}$ 13.2.3, \textquotedblleft topological\textquotedblright{} variant). This completes the proof of (3.2.1).

\sgaStableTarget{sga6.e02.para.3-2-2}3.2.2. For a topological space $S$, we shall denote by $\Parf(S)$ the triangulated category of perfect complexes of $(\underline{O}_S\otimes\underline{\underline{C}})$-Modules, and by $K^\bullet(S)$ the corresponding \textquotedblleft Grothendieck group\textquotedblright{}\SGAAuthorFootnote{1}{see (I 6.3). For a general review of the functors $K$, see Exp.~IV.}.

Let $S$, $X$, $E^\bullet$ be as in (3.2.1); the \uline{analytic index} of $E^\bullet$ is the element
\[
  \operatorname{ind.an}(E^\bullet)=C\ell(\underline{\underline{R}}p_*(E^\bullet))\in K^\bullet(S).
\]
%%%%%%%%%%  scan idx 231  |  folio 218  %%%%%%%%%%
\SGAPrintPage{218}%
For \uline{compact} $S$, $\operatorname{ind.an}(E^\bullet)$ canonically defines an element in the usual group $K(S)$ of topologists. Indeed, in general, for every topological space $S$, denote by $\mathrm{Loclib}(S)$ the additive category of locally free $(\underline{O}_S\otimes\underline{\underline{C}})$-Modules of finite type, and by $K(S)$ the \textquotedblleft Grothendieck group\textquotedblright{} of the triangulated category $K^b(\mathrm{Loclib}(S))/K^{b,\varnothing}(\mathrm{Loclib}(S))$ (the latter is, moreover, identified with $K^b(\mathrm{Loclib}(S))$ when $S$ is paracompact (II 2.3.2 (a)), so that in this case (I 6.4) $K(S)$ is canonically isomorphic to the \textquotedblleft Grothendieck group\textquotedblright{} of the additive category $\mathrm{Loclib}(S)$). There is a natural homomorphism
\[
  K(S)\longrightarrow K^\bullet(S).
\]
But, by (II 2.3.2), this is an isomorphism for compact $S$, proving our assertion. More generally, $\operatorname{ind.an}(E^\bullet)$ defines an element in the completed group
\[
  \widehat K(S)=\varprojlim K(A),
\]
the projective limit being taken over the compact subsets $A$ of $S$.

\sgaStableTarget{sga6.e02.para.3-2-3}3.2.3. The writer does not know whether the finiteness theorem (3.2.1) remains valid when, instead of assuming $S$ locally compact, one assumes only that every point of $S$ has a fundamental system of paracompact neighborhoods. Moreover, the local compactness hypothesis on $S$ does not seem very natural. The space of \textquotedblleft universal local parameters\textquotedblright{} for elliptic complexes of fixed order is indeed an (infinite-dimensional) subspace of a Fréchet space. Let us make this precise. Let $V$ be a compact variety, $F=(F^i)_{i\in\underline{\underline{Z}}}$ a sequence of complex vector bundles of finite rank on $V$, all but finitely many zero, and $\underline n=(n_i)$ a sequence of integers $\geq 0$. Denote by $\operatorname{Comp.diff}^{\underline n}(F)$ the space of complexes of differential operators \textquotedblleft with base $F$\textquotedblright{} and of order $\underline n$, i.e.\ the closed subspace of $\prod\Diff^{n_i}(F^i,F^{i+1})$ defined by the equations $d^{i+1}d^i=0$. (the space $\prod\Diff^{n_i}(F^i,F^{i+1})$ is a Fréchet space). The subset
